\section{Step 11 Preserve the search across a nested call}\label{h263-motion-s01}
Begin with a frame sketch. Mark where you will retain the caller's return address,
the request values, the result pointer, the candidate coordinates and the current
best result. A candidate comparison reuses four argument registers, so the
original request pointer cannot remain live only in \texttt{\$a0}. The SAD entry
also has no obligation to preserve your scratch registers.

Write the prologue and matching epilogue before the loops. Check their saved
register addresses and keep the stack suitably aligned. Return through the saved
caller continuation, rather than through the continuation written by the last
nested call. Do not use an unallocated area below the stack pointer for locals.

\paragraph{Observe} A valid radius-zero request makes one nested comparison.
The function returns to its outer caller with the original saved registers and
stack pointer, even though \texttt{\$ra} changed during the comparison.
\Needspace{8\baselineskip}
\paragraph{Check} Where can a candidate coordinate safely remain across SAD?
\begin{itemize}
\item[A.] In any scratch register, because the caller used it first.
\item[B.] In a preserved register or an allocated stack local.
\item[C.] Only in an input pixel byte.
\end{itemize}

\section{Step 12 Admit the complete request before pixels}\label{h263-motion-s02}
Check the two record pointers, then load the eight request fields. Keep invalid
status as the provisional return. Validate the fields and both image extents
before entering a candidate loop. A request with a valid current window but an
invalid reference stride must fail without reading either image.

Use unsigned range checks for unsigned fields. Values such as
$2^{32}+2$ are not radius two. Address additions must detect carry out of the
64-bit address range; merely using a doubleword instruction does not detect it.
Leave the output record unchanged on every rejection path.

\paragraph{Observe} Radius three, reference stride 175 and an unaligned result
pointer each return status one, candidate count zero and the original result bytes.
\Needspace{8\baselineskip}
\paragraph{Check} Which read is permitted while rejecting a bad mode field?
\begin{itemize}
\item[A.] The mapped request field needed to inspect that mode.
\item[B.] The first reference vector, to see whether the image looks usable.
\item[C.] An arbitrary address beyond the reference plane.
\end{itemize}

\section{Step 13 Clip signed displacement intervals}\label{h263-motion-s03}
Compute the four interval bounds from the validated coordinates and radius.
The lower bounds may be negative, so store and compare them as signed values.
The picture coordinates obtained after adding a displacement must remain within
the valid top-left ranges. Keep the right and bottom limits at 160 and 128,
which already account for the sixteen-sample block size.

Test a corner before an interior block. At the upper-left corner and radius two,
negative displacements disappear. At the lower-right corner, positive ones
disappear. A radius-one interior block has nine candidates; it is a different
geometry from the radius-two corner even though the count agrees.

\paragraph{Observe} Counts are 25 for an interior radius-two search, 15 at one
straight edge, nine at a corner and one for radius zero.
\Needspace{8\baselineskip}
\paragraph{Check} Why is a top-left reference coordinate of 175 invalid?
\begin{itemize}
\item[A.] Byte 175 can never be mapped.
\item[B.] Every motion vector must be zero.
\item[C.] A sixteen-sample block starting there extends past the picture width.
\end{itemize}

\section{Step 14 Compare the zero candidate once}\label{h263-motion-s04}
Initialize the best score above the SAD maximum, then evaluate zero displacement.
Build its reference pointer from the plane origin and the current picture
coordinates. Supply the current pointer unchanged. Pass each stride separately
to the selected comparison mode. Retain the score and zero displacement as the
first completed best result.

If your later rectangle traversal includes $(0,0)$, skip that one entry because
it has already been considered. The successful candidate count describes unique
admitted positions. It does not include an extra visit for initialization.
Equivalent implementations may compute this count from the interval lengths.

\paragraph{Observe} The worked zero candidate has score 7872. Flat equal images
have score zero and ultimately retain displacement $(0,0)$.
\Needspace{8\baselineskip}
\paragraph{Check} What is a safe initial best-score sentinel?
\begin{itemize}
\item[A.] Zero, because SAD is nonnegative.
\item[B.] 65\,536, which exceeds the maximum possible block SAD.
\item[C.] The current image address.
\end{itemize}

\section{Step 15 Form each candidate from the plane origin}\label{h263-motion-s05}
For every remaining admitted $(d_x,d_y)$, form $(x+d_x,y+d_y)$ first, then its
byte offset using the reference stride. Add that offset to the reference plane
pointer. Recomputing from the origin avoids carrying a pointer from the end of
one candidate's sixteen rows into the next candidate.

For the worked position $(16,32)$ and displacement $(-2,-2)$, the candidate begins
at $(14,30)$. With stride 176, its offset is $30\times176+14=5294$ bytes from the
reference origin. Its last row begins another $15\times176=2640$ bytes later.
The current pointer continues to identify the original current block.

\paragraph{Observe} A load from the candidate's first row reads reference
samples 14 through 29 of row 30. It reads no column from the following row.
\Needspace{8\baselineskip}
\paragraph{Check} Which pointer changes when selecting a new candidate?
\begin{itemize}
\item[A.] The reference window pointer derived from the plane origin.
\item[B.] Both image pointers by the same displacement.
\item[C.] The result pointer by sixteen rows.
\end{itemize}

\section{Step 16 Update the best result without changing ties}\label{h263-motion-s06}
Compare a completed SAD with the current best. In centre-first, increasing-row
order, replace only for a strictly smaller value. Copy both displacement
components whenever you replace the score. Updating the score alone leaves a
plausible number attached to the wrong reference block.

Check a periodic pattern as well as a flat one. A checker pattern and its
opposite can have several equally good nonzero displacements while the zero
candidate is worse. The correct choice then uses increasing vertical and
horizontal displacement. A reverse traversal needs an explicit equal-score
comparison to produce that same choice.

\paragraph{Observe} With the public opposite checker patterns at $(16,32)$,
radius two selects $(-1,-2)$ with score zero. Flat equal inputs select $(0,0)$.
\Needspace{8\baselineskip}
\paragraph{Check} What happens if equal scores always replace the best?
\begin{itemize}
\item[A.] SAD becomes negative.
\item[B.] The vector registers become read-only.
\item[C.] Traversal order can incorrectly replace the specified preferred tie.
\end{itemize}

\section{Step 17 Finish the rectangle without crossing an edge}\label{h263-motion-s07}
Advance the horizontal coordinate through its inclusive interval. At its end,
reset it to the first horizontal bound and advance the vertical coordinate.
Stop after the final admitted row. A signed negative loop coordinate needs a
signed comparison; interpreting it as an unsigned address-sized integer changes
the ordering.

A zero score is already minimal, but an early exit still has to preserve the
tie rule and return the full unique candidate count. For this first implementation,
finish the small rectangle. Later work can reduce comparisons while retaining
the complete observable contract.

\paragraph{Observe} A radius-zero request neither repeats its centre nor enters
a nonexistent next row. The corner cases never read outside the clipped union.
\Needspace{8\baselineskip}
\paragraph{Check} What does the returned candidate count describe?
\begin{itemize}
\item[A.] The number of scratch registers used.
\item[B.] The number of unique positions admitted by the clipped search.
\item[C.] The number of output bytes.
\end{itemize}

\section{Step 18 Publish one complete result}\label{h263-motion-s08}
Keep the provisional best state private until the search succeeds. Then store
both signed displacement components and the SAD into the three result fields.
Publish success and the candidate count in the return registers, restore the
outer caller's state and return. Negative displacements use full-width signed
representations; a 32-bit zero-extended negative value is not the required
64-bit result.

The result is a candidate decision, not an already coded motion-vector
difference. Prediction from neighbouring macroblocks, vector units and variable
length coding remain separate later responsibilities. Keeping these boundaries
visible makes the first result inspectable before a bitstream exists.

\paragraph{Observe} The worked caller's result record contains $-2$, $-2$ and
2976. Its returned count is 25 in both scalar and MSA modes.
\Needspace{8\baselineskip}
\paragraph{Check} What may an invalid request publish in the result area?
\begin{itemize}
\item[A.] Nothing; every existing result byte must remain unchanged.
\item[B.] A partially updated displacement.
\item[C.] A success score with an invalid status.
\end{itemize}

\section{Step 19 Run independent search feedback}\label{h263-motion-s09}
Run the completed-entry checks against your working source:
\begin{verbatim}
python3 construction/h263-motion-construction/grade.py \
    my-search.asm --cli /path/to/mars-cli.jar \
    --report search-results.json
\end{verbatim}
The 56 public cases cover both comparison modes and both byte orders, all four
movement quadrants, clipped corners, radius zero, maximal differences, two kinds
of ties and rejected requests. The checker computes every candidate score from
its own input bytes. It checks the selected result, candidate count, input and
request preservation, permitted pixel accesses, saved registers and result guards.
A new report path is required.

Use failures to identify the boundary that differs. A correct SAD with the wrong
sign can still locate the wrong reference block. A correct interior search can
still overread a corner. An implementation may choose its own traversal and
frame layout; it must satisfy these observable conditions.

\paragraph{Observe} Development status four fails completed feedback. A reverse
traversal with the correct explicit tie key can pass the same cases.
\Needspace{8\baselineskip}
\paragraph{Check} What makes a different traversal an equivalent solution?
\begin{itemize}
\item[A.] It uses the reference implementation's private labels.
\item[B.] It has the same instruction count.
\item[C.] It returns the specified result and preserves the same boundaries.
\end{itemize}

\section{Step 20 Measure search work at its own boundary}\label{h263-motion-s10}
Count instructions from search entry through the return to the ordinary caller.
This includes request validation, candidate setup, nested SAD work, selection
and restoration. Compare the same input and radius in scalar and MSA modes.
Do not multiply the earlier kernel count and report that product as a measured
whole-search count: it omits the work that surrounds each comparison.

The public feedback reports completed instruction counts for its named inputs.
A small caller with radius zero exposes fixed search overhead; an interior
radius-two case adds candidate work; a corner exposes clipping. Retain correctness
and access checks while comparing these results. The count measures simulator
instruction work, not host cycles, a video frame rate or completed encoder time.

\paragraph{Observe} Record the count and candidate geometry together. Replacing
SAD with a faster correct MSA body can reduce comparison work while leaving the
same displacement, score and candidate count.
\Needspace{8\baselineskip}
\paragraph{Check} Which measurement includes search setup and restoration?
\begin{itemize}
\item[A.] The file size of the generated raw fixture.
\item[B.] A trace covering search entry through return to its caller.
\item[C.] The SAD kernel count alone, without entering search.
\end{itemize}
